\section*{第 \ref{ch:b1:poly} 章 --- 多项式}

\begin{solution}{exo:b1:poly:1}
$X^5 - 1 = (X^2 + X + 1)(X^3 - X^2 + 1) + (-X - 2)$。步骤：减去
$X^3 B$，再减去 $-X^2 B$，最后减去 $B$；余式 $-X - 2$ 的次数为
$1 < 2$。\emph{在 $X = 1$ 时验算：}$\;0 = 3 \times 1 + (-3)$。

$2X^4 + X^3 - X + 3 = (X^2 - 2)(2X^2 + X + 4) + (X + 11)$。\emph{在
$X = 0$ 处验算：} $\;3 = (-2)(4) + 11$。
\end{solution}

\begin{solution}{exo:b1:poly:2}
$X^2 + X + 1 = (X - j)(X - j^2)$，其中 $j = \eu^{2\iu\pi/3}$，$j^3 =
1$。它整除 $Q_n = X^{2n} + X^n + 1$，当且仅当 $j$ 与 $j^2$ 都是
$Q_n$ 的根；因 $Q_n$ 的系数为实数，$Q_n(j^2) =
\conj{Q_n(j)}$，故只需 $Q_n(j) = 0$。现 $Q_n(j) =
j^{2n} + j^n + 1$ 取决于 $n$ 模 $3$ 的余数：
\begin{itemize}
  \item $n \equiv 0$: $Q_n(j) = 1 + 1 + 1 = 3 \neq 0$;
  \item $n \equiv 1$: $Q_n(j) = j^2 + j + 1 = 0$;
  \item $n \equiv 2$: $Q_n(j) = j^4 + j^2 + 1 = j + j^2 + 1 = 0$.
\end{itemize}
于是 $X^2 + X + 1 \mid X^{2n} + X^n + 1$ 当且仅当 $3 \nmid n$。
\end{solution}

\begin{solution}{exo:b1:poly:3}
由\cref{prop:b1:poly:multiplicity}可知，$(X-1)^2 \mid P$ 当且仅当 $P(1) =
P'(1) = 0$：
\[
P(1) = 2 + a + b = 0, \qquad P'(1) = 4 + 3a + 2b = 0 .
\]
解得 $b = -a - 2$，并代入得 $4 + 3a - 2a - 4 = a = 0$，故 $a = 0$，$b =
-2$，于是 $P = X^4 - 2X^2 + 1 = (X^2 - 1)^2 = (X-1)^2 (X+1)^2$，这是
实数范围内的因式分解。
\end{solution}

\begin{solution}{exo:b1:poly:4}
$X^3 - 1 = (X - 1)(X - j)(X - j^2)$ 在 $\C$ 上成立（$j = \eu^{2\iu\pi/3}$），
而 $(X - 1)(X^2 + X + 1)$ 是在 $\R$ 上的分解。

$X^4 + X^2 + 1 = (X^2 + X + 1)(X^2 - X + 1)$ 在 $\R$ 上成立（展开即可，
或注意 $X^4 + X^2 + 1 = (X^2+1)^2 - X^2$）；在 $\C$ 上，每个二次式
都可分解：根为 $j, j^2$ 与 $-j, -j^2$，即 $\eu^{\pm 2\iu\pi/3},
\eu^{\pm\iu\pi/3}$。

$X^6 - 1 = \prod_{k=0}^{5} (X - \eu^{\iu k\pi/3})$ 在 $\C$ 上成立，而在
$\R$ 上：
\[
X^6 - 1 = (X-1)(X+1)(X^2 + X + 1)(X^2 - X + 1),
\]
将共轭对 $\eu^{\pm 2\iu\pi/3}$ 与 $\eu^{\pm
\iu\pi/3}$ 分组合并。
\end{solution}

\begin{solution}{exo:b1:poly:5}
\begin{enumerate}
  \item 设 $p/q$（既约分数）为首一整系数多项式
        $X^3 + \dots + c_0$ 的根：
        $P(p/q) = 0$ 中去分母，得 $p^3 = -q\,(\text{整数})$，故 $q \mid
        p^3$；由互素性得 $q = \pm 1$：此根为整数
        $p$，且 $p \mid c_0$（把 $c_0$ 单独移到一边）。本题中候选根为
        $6$ 的因数：逐一检验，$P(1) = 0$，$P(2) = 0$，$P(3) = 0$。于是
        $P = (X-1)(X-2)(X-3)$。
  \item 由韦达定理：$s_1 = 6$，$s_2 = 11$，$s_3 = 6$。平方和：
        $s_1^2 - 2s_2 = 36 - 22 = 14 = 1 + 4 + 9$，与预期相符。倒数之和：
        $\frac{s_2}{s_3} = \frac{11}{6} = 1 + \frac12 + \frac13$，与预期相符。
\end{enumerate}
\end{solution}

\begin{solution}{exo:b1:poly:6}
欧几里得算法的第一步除法：
\[
(X + 1)(X^3 - X^2 + X - 1)
= X^4 - X^3 + X^2 - X + X^3 - X^2 + X - 1 = X^4 - 1 ,
\]
于是 $X^4 - 1$ 除以 $X^3 - X^2 + X - 1$ 恰好除尽
（商为 $X + 1$，余式为 $0$），算法立即停止：
\[
\gcd(X^4 - 1,\; X^3 - X^2 + X - 1) = X^3 - X^2 + X - 1
\]
（已为首一）。贝祖等式是平凡的：$\gcd = 0
\cdot (X^4 - 1) + 1 \cdot (X^3 - X^2 + X - 1)$。以因式分解作一致性检验：
$X^3 - X^2 + X - 1 = (X - 1)(X^2 + 1)$，它确实是
$X^4 - 1 =
(X-1)(X+1)(X^2+1)$ 的公共不可约因式之积。
\end{solution}

\begin{solution}{exo:b1:poly:7}
由于 $P \geq 0$ 在 $\R$ 上成立，其实根的重数均为偶数（在
奇数重根处 $P$ 变号）。利用
\cref{cor:b1:poly:factorization} 并两两配对，写作
\[
P = c \prod_i (X - a_i)^{2k_i} \prod_j \bigl((X - z_j)(X - \conj
z_j)\bigr)^{n_j},
\]
其中 $c > 0$（由 $+\infty$ 处的性态）。令
\[
S = \sqrt c\, \prod_i (X - a_i)^{k_i} \prod_j (X - z_j)^{n_j}
\in \C[X],
\]
从而 $P = S\,\conj S$，其中 $\conj S$ 的系数
多项式。将 $S = A + \iu B$（其中 $A, B \in \R[X]$），则
\[
P = (A + \iu B)(A - \iu B) = A^2 + B^2 .
\]
\end{solution}

\begin{solution}{exo:b1:poly:8}
由拉格朗日插值公式（\cref{thm:b1:poly:lagrange}），取节点 $0, 1, 2$：
\[
P = 1\cdot\frac{(X-1)(X-2)}{(0-1)(0-2)} + 3\cdot\frac{X(X-2)}{1\cdot(1-2)}
+ 2\cdot\frac{X(X-1)}{2\cdot 1}
= \frac{(X-1)(X-2)}{2} - 3X(X-2) + X(X-1).
\]
展开得 $\frac{X^2 - 3X + 2}{2} - 3X^2 + 6X + X^2 - X =
-\frac{3}{2}X^2 + \frac{7}{2}X + 1$.

方程组：$P = aX^2 + bX + c$，其中 $c = 1$；$a + b + 1 = 3$；$4a + 2b +
1 = 2$。第三式减去第二式的两倍，得 $2a - 1 = -4$，
故 $a = -\frac32$，$b = \frac72$，与前面所得多项式相同：
$P = -\frac32 X^2 + \frac72 X + 1$。（验算：$P(2) = -6 + 7 + 1 = 2$。）
\end{solution}

\begin{solution}{exo:b1:poly:9}
$P = X^{2n+1} - 1$：$P' = (2n+1)X^{2n} \geq 0$，故该多项式
函数单调递增（除 $0$ 外严格递增），极限为
$\mp\infty$，故它在 $\R$ 上恰有一个零点（$x = 1$）。

设 $E_n = \sum_{k=0}^{n} \frac{X^k}{k!}$，则 $E_n' = E_{n-1} = E_n
- \frac{X^n}{n!}$。若 $a$ 为重根，则 $E_n(a) =
E_n'(a) = 0$（见\cref{prop:b1:poly:multiplicity}），于是
$\frac{a^n}{n!} = E_n(a) - E_n'(a) = 0$，故 $a = 0$；但 $E_n(0) = 1
\neq 0$。故无重根。
\end{solution}

\begin{solution}{exo:b1:poly:10}
\begin{enumerate}
  \item 归纳法（两个初始情形成立）。利用 $\cos(n+1)\theta +
        \cos(n-1)\theta = 2\cos\theta\cos n\theta$:
        \[
        T_{n+1}(\cos\theta) = 2\cos\theta \cos n\theta -
        \cos(n-1)\theta = \cos(n+1)\theta .
        \]
  \item $T_n(\cos\theta) = 0$ 当且仅当 $n\theta \equiv \frac\pi2 \pmod
        \pi$：这些数
        \[
        x_k = \cos\Bigl(\frac{(2k+1)\pi}{2n}\Bigr),
        \qquad k = 0, 1, \dots, n-1,
        \]
        为 $n$ 个互异点，且均属于 $\intoo{-1}{1}$（角度位于
        $\intoo{0}{\pi}$ 内，此时 $\cos$ 为单射），它们都是
        $T_n$ 的根。由于 $\deg T_n = n$（由递推式得出，
        其首项系数为 $2^{n-1}$，$n \geq 1$ 时成立，由归纳法可得），
        这些正是\emph{全部}根，且均为单根。
  \item 对 $x = \cos\theta \in \intcc{-1}{1}$，有 $\abs{T_n(x)} =
        \abs{\cos n\theta} \leq 1$，等号成立当且仅当 $n\theta \equiv
        0 \pmod\pi$，即在 $n+1$ 个点 $y_k =
        \cos\frac{k\pi}{n}$，$k = 0, \dots, n$，此时 $T_n(y_k) =
        (-1)^k$。（这一等幅振荡性质使 $2^{1-n}T_n$ 成为
        首一的 $n$ 次多项式，它在
        $\intcc{-1}{1}$ 上确界范数最小，
        此事实已在本章周末问题中证明。）
\end{enumerate}
\end{solution}

\begin{solution}{exo:b1:poly:11}
记 $P = c\prod_i (X - a_i)^{m_i}$。由乘积求导法则（推广至
多个因子）得
\[
P' = c\sum_{i} m_i (X - a_i)^{m_i - 1} \prod_{k \neq i} (X -
a_k)^{m_k},
\]
两边除以 $P$，得 $\frac{P'}{P} = \sum_i \frac{m_i}{X - a_i}$
（作为有理函数，即在非根处）。

设 $w$ 为 $P'$ 的根。若 $w$ 为某个 $a_i$，则它显然位于
凸包中。否则，在 $w$ 处取值得
\[
0 = \sum_i \frac{m_i}{w - a_i}
= \sum_i m_i\, \frac{\conj w - \conj a_i}{\abs{w - a_i}^2} .
\]
取共轭得 $\sum_i \lambda_i (w - a_i) = 0$，其中 $\lambda_i =
\frac{m_i}{\abs{w - a_i}^2} > 0$。于是
\[
w = \frac{\sum_i \lambda_i a_i}{\sum_i \lambda_i} :
\]
该点为凸组合（正权重经归一化后之和为 $1$），
其组合元素为根 $a_i$。因此 $P'$ 的每个根均位于
$P$ 的根所构成的凸包中。
\end{solution}

\begin{solution}{exo:b1:poly:12}
将 $(1 + X)^n$ 在三个三次
单位根处的值相加：
\[
2^n + (1 + j)^n + (1 + j^2)^n
= \sum_{k=0}^n \binom nk\,\bigl(1 + j^k + j^{2k}\bigr)
= 3\sum_{k\,:\,3\mid k}\binom nk ,
\]
由于 $1 + j^k + j^{2k}$ 是等比和，其值为 $3$，当 $3
\mid k$ 时成立；否则为 $\frac{j^{3k} - 1}{j^k - 1} = 0$。又 $1 +
j = \frac12 + \iu\frac{\sqrt3}2 = \eu^{\iu\pi/3}$，且 $1 + j^2 =
\conj{1 + j} = \eu^{-\iu\pi/3}$，故 $(1+j)^n + (1+j^2)^n =
2\cos\frac{n\pi}3$，于是
\[
\sum_{k\geq0}\binom n{3k} = \frac{2^n + 2\cos\frac{n\pi}3}{3} .
\]
验算：$n = 3$ 时，$\frac{8 + 2\cos\pi}3 = 2 = \binom30 + \binom33$；
$n = 6$ 时，$\frac{64 + 2}3 = 22 = 1 + 20 + 1$。
\end{solution}

\begin{solution}{pb:b1:poly:1}
\textbf{1.} $T_2 = 2X^2 - 1$; $T_3 = 2X(2X^2 - 1) - X = 4X^3 -
3X$; $T_4 = 2X\,T_3 - T_2 = 8X^4 - 8X^2 + 1$; $T_5 = 2X\,T_4 - T_3
= 16X^5 - 20X^3 + 5X$。恒等式 $T_3(\cos\theta) = \cos3\theta$
恰为 $\cos3\theta = 4\cos^3\theta - 3\cos\theta$，
即 \cref{ex:b1:complex:cos3} 中的结论。

\textbf{2.} 当 $n = 1, 2$ 时成立。若 $T_{n-1}$、$T_n$ 的次数分别为
$n-1$、$n$，首项系数分别为 $2^{n-2}$、$2^{n-1}$，则
$2X\,T_n$ 的次数为 $n+1$，首项系数为 $2^n$，而
$T_{n-1}$ 次数更低；故 $T_{n+1}$ 次数为 $n + 1$，首项
系数为 $2^n$。关于奇偶性：若 $T_{n-1}$ 与 $n - 1$ 同奇偶，
且 $T_n$ 与 $n$ 同奇偶，则 $2X\,T_n$ 与 $T_{n-1}$ 均与
$n + 1$ 同奇偶，因此 $T_{n+1}$ 亦然。

\textbf{3.} 若 $P(\cos\theta) = \cos n\theta$ 对所有 $\theta$ 成立，
则 $P$ 与 $T_n$ 在 $\intcc{-1}1$ 的每一点处取值相同 --- 该集合
为无限集 --- 故 $P - T_n$ 有无穷多个根，
因而为零多项式（见\cref{cor:b1:poly:nroots}）。

\textbf{4.} 设 $x = \cos\theta$，则 $T_m(T_n(\cos\theta)) =
T_m(\cos n\theta) = \cos(mn\theta) = T_{mn}(\cos\theta)$，且
$2T_mT_n(\cos\theta) = 2\cos m\theta\cos n\theta = \cos(m+n)\theta
+ \cos\abs{m - n}\theta$。两式在 $\intcc{-1}1$ 上均成立，
故由第 3 题的论证，二者亦为多项式恒等式。

\textbf{5.} $x_k$ 为 $n$ 个互异的单根，且
首项系数为 $2^{n-1}$：
\[
T_n = 2^{n-1}\prod_{k=0}^{n-1}
\Bigl(X - \cos\frac{(2k+1)\pi}{2n}\Bigr) .
\]
交错性：诸角 $0 < \frac{\pi}{2n} < \frac\pi n <
\frac{3\pi}{2n} < \frac{2\pi}n < \dots < \pi$ 中的角交替地属于
$y$ 角 $\frac{k\pi}n$ 与 $x$ 角
$\frac{(2k+1)\pi}{2n}$；由于 $\cos$ 在
$\intcc0\pi$ 上严格递减，故诸值按相反次序交错：$y_n <
x_{n-1} < y_{n-1} < \dots < x_0 < y_0$。相邻两个极值点之间
恰有一个根，这与 $\cos n\theta$ 的图像
所示相符。

\textbf{6.} 用归纳法，并借助 $2\cosh a\cosh b = \cosh(a + b) +
\cosh(a - b)$（见\cref{prop:b1:functions:hyprules}）：$T_{n+1}(\cosh
t) = 2\cosh t\cosh nt - \cosh(n-1)t = \cosh(n+1)t$。对 $x \geq
1$，记 $x = \cosh t$，其中 $t \geq 0$；则 $\eu^t = x +
\sqrt{x^2 - 1}$，又 $\eu^{-t} = x - \sqrt{x^2 - 1}$，于是
\[
T_n(x) = \cosh(nt)
= \frac{(x + \sqrt{x^2-1})^n + (x - \sqrt{x^2-1})^n}2 .
\]
当 $x > 1$ 时，第一项严格大于 $\frac12(1)^n$，且
按几何级数增长：$T_n(x) > 1$。

\textbf{7.} 棣莫弗公式：$\cos n\theta = \Re\bigl((\cos\theta +
\iu\sin\theta)^n\bigr) = \sum_{2j \leq n}\binom n{2j}
\cos^{n-2j}\theta\,(\iu\sin\theta)^{2j}$，又 $(\iu\sin\theta)^{2j}
= (-\sin^2\theta)^j = (\cos^2\theta - 1)^j$。代入 $x =
\cos\theta$，并运用第 3 题的结论：
\[
T_n(x) = \sum_{0\leq 2j\leq n}\binom n{2j}x^{n-2j}(x^2 - 1)^j .
\]
当 $n = 3$ 时：$\binom30 x^3 + \binom32 x(x^2 - 1) = x^3 + 3x^3 -
3x = 4x^3 - 3x$，与第 1 题一致。

\textbf{8.} $T_n(1) = \cos(n\cdot0) = 1$; $T_n(-1) = \cos(n\pi) =
(-1)^n$；$T_n(0) = \cos\frac{n\pi}2$，其值为 $0$（当 $n$ 为奇数时），而
为 $(-1)^{n/2}$（当 $n$ 为偶数时）。

\textbf{9.} 用归纳法证明 $U_n(\cos\theta) =
\frac{\sin(n+1)\theta}{\sin\theta}$：对 $U_0 = 1$ 与 $U_1 =
2X$（$\sin2\theta = 2\sin\theta\cos\theta$）成立；归纳步骤为
和差化积公式 $\sin(n+2)\theta = 2\cos\theta\,
\sin(n+1)\theta - \sin n\theta$。现将
$T_n(\cos\theta) = \cos n\theta$ 关于 $\theta$ 求导：
$-\sin\theta\,T_n'(\cos\theta) = -n\sin n\theta$，于是当 $\theta
\notin \pi\Z$ 时：
\[
T_n'(\cos\theta) = n\,\frac{\sin n\theta}{\sin\theta}
= n\,U_{n-1}(\cos\theta) ,
\]
又多项式 $T_n'$ 与 $nU_{n-1}$ 在
$\intoo{-1}1$ 上处处相等，故二者恒等。

\textbf{10.} $\abs{\sin(n+1)\theta} = \abs{\sin n\theta\cos\theta
+ \cos n\theta\sin\theta} \leq \abs{\sin n\theta} +
\abs{\sin\theta}$，由归纳法得 $\abs{\sin n\theta} \leq
n\abs{\sin\theta}$。因此 $\abs{U_{n-1}} \leq n$ 在 $\intoo{-1}1$ 上
成立，且 $\abs{T_n'} = n\abs{U_{n-1}} \leq n^2$ 在该区间上亦然；在 $\pm1$ 处，
此界可由取极限推出（或直接：由递推式得 $U_{n-1}(1) = n$，
由归纳法得 $U_n(1) = n + 1$，而由奇偶性得
$U_{n-1}(-1) = (-1)^{n-1}n$）。于是 $T_n'(1) = n^2$，$T_n'(-1)
= (-1)^{n-1}n^2$：界 $n^2$ 在端点处取到。

\textbf{11.} 将 $\sin\theta\,T_n'(\cos\theta) = n\sin
n\theta$（第 9 题）关于 $\theta$ 求导：
\[
\cos\theta\,T_n'(\cos\theta) - \sin^2\theta\,T_n''(\cos\theta)
= n^2\cos n\theta = n^2\,T_n(\cos\theta) .
\]
取 $x = \cos\theta$，并利用 $\sin^2\theta = 1 - x^2$，得 $x\,T_n' - (1
- x^2)T_n'' = n^2T_n$ 在 $\intcc{-1}1$ 上成立，从而处处成立：
$(1 - x^2)y'' - xy' + n^2y = 0$，其中 $y = T_n$。验证 $T_2 =
2x^2 - 1$: $(1 - x^2)(4) - x(4x) + 4(2x^2 - 1) = 4 - 4x^2 - 4x^2
+ 8x^2 - 4 = 0$.

\textbf{12.} $\widetilde T_n$ 为首一多项式（第 2 题），且
$\abs{\widetilde T_n} = 2^{1-n}\abs{T_n} \leq 2^{1-n}$ 在
$\intcc{-1}1$ 上，且 $\widetilde T_n(y_k) = (-1)^k2^{1-n}$ 在
$n + 1$ 个点 $y_k$（见\cref{exo:b1:poly:10}）上取值，故范数恰为
$2^{1-n}$，并在符号正负交替的点上取到。

\textbf{13.} $\widetilde T_n$ 与 $P$ 均为首一的
$n$ 次多项式，故首项相消，$\deg D \leq n - 1$。在 $y_k$ 处：
$D(y_k) = (-1)^k2^{1-n} - P(y_k)$，又 $\abs{P(y_k)} \leq \norm
P_\infty < 2^{1-n}$ 迫使 $D(y_k)$ 的符号与
$(-1)^k2^{1-n}$ 的符号相同，且为严格的。

\textbf{14.} $D$ 在 $y_{k+1}$ 与 $y_k$ 之间变号，对每个
$k = 0, \dots, n-1$：由介值性质，$D$ 在
这 $n$ 个两两不交的开区间中各有一根 --- $n$
个互异的根，而次数 $\leq n - 1$ 的非零多项式
不可能。又 $D = 0$ 亦不可能（两者范数不同）。
矛盾。故任一首一的 $P$（次数为 $n$）均不满足 $\norm P_\infty <
2^{1-n}$，这正是切比雪夫定理。

\textbf{15.} 现仅知 $\abs{P(y_k)} \leq 2^{1-n}$，故 $(-1)^k
D(y_k) = 2^{1-n} - (-1)^kP(y_k) \geq 2^{1-n} - \abs{P(y_k)} \geq
0$。若 $D(y_k) = 0$，且 $y_k$ 为内点（$0 < k < n$），则
$P(y_k) = (-1)^k2^{1-n}$，故 $\abs P$ 的上确界
$2^{1-n}$ 在内点 $y_k$ 处取到，因而 $P'(y_k) = 0$
（内点极值）；又 $T_n'(y_k) = nU_{n-1}(y_k) = 0$，因为
$\sin(n\cdot\frac{k\pi}n) = 0$ --- 故 $\widetilde T_n'(y_k) = 0$
亦然，从而 $D'(y_k) = 0$：$y_k$ 为 $D$ 的根，其重数至少为
$2$。

\textbf{16.} 计算 $D$ 的根（计重数）。设 $z$ 为
内点 $y_k$ 中满足 $D(y_k) = 0$ 者的个数（每个都是
二重根，由第 15 题），$e \in \{0, 1, 2\}$ 为
端点（$y_0$ 或 $y_n$）中满足 $D = 0$ 者的个数（每个至少是
单根）。若区间 $(y_{k+1}, y_k)$ 的两个端点处 $D
\neq 0$，则两端的符号严格交错，因而区间内部必有一根。
每个为零的内点至多破坏其两侧的两个区间，
每个为零的端点至多破坏一个区间。于是至少
$n - 2z - e$ 个区间各仍贡献一个根，这些根
均与 $y$ 处的根互异。合计至少 $(n - 2z - e) + 2z + e = n$
个根（计重数），而次数 $\leq n - 1$ 的多项式
只能恒为零，故 $D = 0$，$P = \widetilde T_n$。极小化元唯一。

\textbf{17.} 仿射映射 $t \mapsto x = \frac{a+b}2 +
\frac{b-a}2\,t$ 是 $\intcc{-1}1 \to \intcc ab$ 的双射。若 $P$
首一且次数为 $n$，则 $Q(t) = P(x(t))$ 是关于
$t$ 的多项式，首项系数为 $\bigl(\frac{b-a}2\bigr)^n$，且
$\sup_{\intcc ab}\abs P = \sup_{\intcc{-1}1}\abs Q$。首一
多项式 $Q/\bigl(\frac{b-a}2\bigr)^n$ 的上确界范数 $\geq
2^{1-n}$（第 13--14 题），故
\[
\sup_{\intcc ab}\abs P \geq \Bigl(\frac{b-a}2\Bigr)^n 2^{1-n}
= 2\Bigl(\frac{b-a}4\Bigr)^n ,
\]
等号恰在 $P(x) = \bigl(\frac{b-a}2\bigr)^n
\widetilde T_n\bigl(t(x)\bigr)$ 时取等（由第 16 题）。

\textbf{18.} $\widetilde T_3 = \frac{T_3}4 = X^3 - \frac34X$；
$\widetilde T_3{}' = 3X^2 - \frac34$ 在 $\pm\frac12$ 处为零。
取值：$\widetilde T_3(-1) = -\frac14$，$\widetilde
T_3(-\tfrac12) = \frac14$，$\widetilde T_3(\tfrac12) = -\frac14$，
$\widetilde T_3(1) = \frac14$：四个
绝对值为 $\frac14$ 的交错极值 --- 故 $\norm{\widetilde T_3}_\infty =
\frac14$，由切比雪夫定理，没有首一三次多项式
在 $\intcc{-1}1$ 上的上确界范数比它更小。

\textbf{19.} $\omega$ 首一且次数为 $n + 1$，故
$\norm\omega_\infty \geq 2^{-n}$，由切比雪夫定理（次数
为 $n+1$），等号成立当且仅当 $\omega = \widetilde T_{n+1}
= 2^{-n}T_{n+1}$（由第 16 题），即当且仅当节点
恰为 $n + 1$ 个 $T_{n+1}$ 的根。取切比雪夫节点时，误差因子
$\norm\omega_\infty$ 为 $2^{-n}$ --- 这是可能取到的最小值。

\textbf{20.} $5 \times 36^\circ = 180^\circ$，故 $T_5(c) =
\cos180^\circ = -1$：$16c^5 - 20c^3 + 5c + 1 = 0$。代入 $x =
-1$ 得 $-16 + 20 - 5 + 1 = 0$，展开后可确认
\[
16x^5 - 20x^3 + 5x + 1 = (x + 1)\bigl(4x^2 - 2x - 1\bigr)^2 .
\]
由于 $c = \cos36^\circ \neq -1$，故 $c$ 是 $4x^2 - 2x -
1$ 的根，其两根为 $\frac{1 \pm \sqrt5}4$；因 $c > 0$，
\[
\cos36^\circ = \frac{1 + \sqrt5}4 .
\]
检验：$\cos72^\circ = T_2(c) = 2c^2 - 1 = 2\cdot\frac{3 +
\sqrt5}8 - 1 = \frac{\sqrt5 - 1}4$，与
\cref{exo:b1:complex:8}中求得的值一致。

\textbf{21.} $\sqrt{1.1^2 - 1} = \sqrt{0.21} \approx 0.458$，故
$x + \sqrt{x^2-1} \approx 1.558$，且 $(1.558)^{10} \approx 84.5$，
而 $(1.1 - 0.458)^{10} \approx 0.01$：$T_{10}(1.1) \approx
\frac{84.5 + 0.01}2 \approx 42$。一个被限制在
$\intcc{-1}1$ 上有界的多项式，在其边界外
十分之一处已超过 $40$：线段上的有界性
对线段之外稍远之处毫无说明。

\textbf{22.} 在第 7 题给出的 $T_p$ 公式中，$j = 0$ 项
为 $X^p$；其余每一项都含 $\binom p{2j}$，且 $0 < 2j <
p$（注意 $2j \neq p$，因 $p$ 为奇数），它被 $p$ 整除，
由\cref{thm:b1:arith:fermat}证明的第一步可知。
因此 $T_p - X^p$ 的每个系数都是 $p$ 的倍数。
验算：$T_3 - X^3 = 3X^3 - 3X = 3(X^3 - X)$；$T_5 - X^5 = 15X^5
- 20X^3 + 5X = 5(3X^5 - 4X^3 + X)$.

\textbf{23.} 由第 17 题，取 $\intcc ab = \intcc01$ 且 $n =
2$：最小偏差为 $2\bigl(\frac14\bigr)^2 = \frac18$，
可由 $\bigl(\frac12\bigr)^2\widetilde T_2(2x - 1) =
\frac14\bigl((2x-1)^2 - \frac12\bigr) = x^2 - x + \frac18$ 取到。
在 $\intcc01$ 上距零最近的首一二次多项式为 $x^2 - x +
\frac18$，其上确界范数为 $\frac18$。

\textbf{24.} (i) 刚性 --- 根的个数多于次数的多项式
必为零 --- 它支撑了唯一性原理
（第 3 题），三角恒等式向
多项式恒等式的转化（第 4、7、9、11 题），以及
极值性证明中的两个根计数论证（第 14、
16 题）。(ii) 在\cref{ch:b1:complex}中，三角学（棣莫弗公式、
和差化积）以及
在\cref{ch:b1:functions}中的双曲函数为该族背后的每一个恒等式提供了依据；
代换 $x = \cos\theta$ 是其间的桥梁。(iii)
在\cref{ch:b1:arith}中的整除性 $p \mid \binom p{2j}$ 把
系数公式化为第 22 题中的同余式。

\textbf{25.} 切比雪夫定理把一个在无穷维族（全体首一多项式）上的优化问题
化为有限的组合论证：
一个优于 $\widetilde T_n$ 的竞争者，将与它相差一个低次多项式，
而该多项式被迫变号
$n$ 次 --- 比其次数所允许的根多一个。
因此，等振荡模式并非一种巧合，而恰恰是
最优性的判据，而取等情形则使
带重数的根计数更为精细。代换 $x =
\cos\theta$ 值得最后一提：它把刚性的
离散的多项式世界搬入了周期性的三角学世界，
在那里 $T_n$ 的根与极值点不过是
$\cos n\theta$ 的规则格点。所借用的两个分析学事实
--- 介值性质（第 14 题；
在\cref{ch:b1:continuity}中证明），以及在
内点极值处导数为零（第 15 题；在\cref{ch:b1:derivative}中证明）
--- 恰是后续章节将要归还的工具，
使全文首尾呼应。
\end{solution}
